\documentclass[12pt,a4paper]{article}

\usepackage[utf8]{inputenc}
\usepackage[T1]{fontenc}
\usepackage[catalan]{babel}
\usepackage{amsmath}
\usepackage{amssymb}
\usepackage{geometry}
\usepackage{enumitem}
\usepackage{titlesec}
\usepackage{parskip}
\usepackage{fancyhdr}
\usepackage{tabularx}
\usepackage{booktabs}
\usepackage{array}
\usepackage{tcolorbox}
\usepackage{tikz}

\geometry{top=2.0cm, bottom=1.5cm, left=1.5cm, right=1.5cm, headheight=15pt, headsep=0.5cm}

\pagestyle{fancy}
\fancyhf{}
\fancyhead[L]{\small\textit{Física · 4t ESO · Escola Cervetó}}
\fancyhead[R]{\small\textit{Cinemàtica}}
\fancyfoot[C]{\small\thepage}
\renewcommand{\headrulewidth}{0.4pt}

\titleformat{\section}[block]{\large\bfseries}{}{0em}{}[\vspace{2pt}\hrule\vspace{4pt}]

\newcounter{exercici}
% Exercici command with grade box on the same line
\newcommand{\exercici}[2]{%
  \stepcounter{exercici}
  \section*{Exercici \theexercici\ \textnormal{(#1 punts)} \hfill \normalfont\large Nota: \framebox[2.5cm]{\rule{0pt}{20pt}} / #1}
  \vspace{-10pt}
  \textit{#2}
  \vspace{5pt}
}

% Grade box - removed as it's now integrated in the header
\newcommand{\gradebox}[1]{}

% Plots filling most of the page with more boxes
\newcommand{\drawplots}[1]{%
  \vfill
  \begin{center}
    \begin{tikzpicture}[xscale=2.1, yscale=2.4]
      % Position Grid (Top)
      \draw[step=0.2, gray!15, ultra thin] (0,6.6) grid (8.5,9.1);
      \draw[->, thick] (0,7.85) -- (8.8,7.85) node[right] {$t$};
      \draw[->, thick] (0,6.6) -- (0,9.4) node[above] {$#1$};
      
      % Velocity Grid (Middle)
      \draw[step=0.2, gray!15, ultra thin] (0,3.3) grid (8.5,5.8);
      \draw[->, thick] (0,4.55) -- (8.8,4.55) node[right] {$t$};
      \draw[->, thick] (0,3.3) -- (0,6.1) node[above] {$v$};

      % Acceleration Grid (Bottom)
      \draw[step=0.2, gray!15, ultra thin] (0,0) grid (8.5,2.5);
      \draw[->, thick] (0,1.25) -- (8.8,1.25) node[right] {$t$};
      \draw[->, thick] (0,0) -- (0,2.8) node[above] {$a$};
    \end{tikzpicture}
  \end{center}
}

\begin{document}

% ================================================================
% PÀGINA 1: INSTRUCCIONS I PUNTUACIÓ
% ================================================================
\begin{center}
  \vspace*{5pt}
  {\Huge\bfseries Examen de Cinemàtica}\\[14pt]
  {\Large Física · 4t ESO \quad·\quad Escola Cervetó \quad·\quad Prof.\ Oriol Farrés}\\[18pt]
  \hrule height 2pt
  \vspace{10pt}
\end{center}

\vspace{10pt}

\noindent
\begin{tabularx}{\textwidth}{|X|}
  \hline
  \rule{0pt}{18pt}\textbf{\large Nom i cognoms:}\\[8pt]
  \hline
\end{tabularx}

\vspace{15pt}

\section*{Instruccions de l'Examen}

\begin{itemize}[leftmargin=1.5em, itemsep=6pt]
  \item \textbf{Constants:} Usa $g = 9{,}81\ \mathrm{m/s^2}$. Es permet usar $\mathbf{\frac{1}{2} g = 4{,}9}$ per simplificar.
  \item \textbf{Càlculs:} Arrodoneix a \textbf{2 decimals} a cada pas del procés.
  \item \textbf{Unitats:} \textbf{Mostra sempre la conversió al SI} de forma explícita.
  \item \textbf{Estructura Obligatòria:} Cal seguir els 4 passos per a cada exercici: (1)~Fórmules, (2)~Dades/SI, (3)~Esquema del problema, (4)~Resolució de cada apartat.
  \item \textbf{Gràfiques:} Usa els eixos proporcionats. Deuen ser qualitativament correctes. \textbf{És obligatori indicar les mètriques i els punts rellevants als eixos}. Dibuixar només la forma de la gràfica sense valors és fonamentalment invàlid i es considerarà com no realitzat.
\end{itemize}

\section*{Taula de Conversió de Longitud}
\noindent
\begin{tabularx}{\textwidth}{|*{7}{>{\centering\arraybackslash}X|}}
  \hline
  \rule{0pt}{15pt}\textbf{km} & \textbf{hm} & \textbf{dam} & \textbf{m} & \textbf{dm} & \textbf{cm} & \textbf{mm} \\[5pt]
  \hline
  \rule{0pt}{20pt} 0 & 2 & 0 & 0 & 0 & 0 & 0 \\[10pt]
  \hline
  \rule{0pt}{20pt} 0 & 7 & 5 & 0 & 0 & 0 & 0 \\[10pt]
  \hline
  \rule{0pt}{20pt} & & & & & & \\[10pt]
  \hline
  \rule{0pt}{20pt} & & & & & & \\[10pt]
  \hline
  \rule{0pt}{20pt} & & & & & & \\[10pt]
  \hline
\end{tabularx}

\vfill

\section*{Graella de Puntuació}
\noindent
\begin{tabularx}{\textwidth}{|X|X|X|X|X||X|}
  \hline
  \rule{0pt}{25pt} \textbf{Ex 1:} ~/~2,5 & \textbf{Ex 2:} ~/~2,5 & \textbf{Ex 3:} ~/~2,5 & \textbf{Ex 4:} ~/~2,5 & \textbf{Extra:} ~/~2 & \textbf{TOTAL:} ~/~10 \\[10pt]
  \hline
\end{tabularx}

\newpage

% ================================================================
% PÀGINA 2: EXERCICI 1
% ================================================================
\exercici{2,5}{MRU · Einstein vs Bohr (Conferència de Solvay)}

L'octubre de 1927, durant la cèlebre Cinquena Conferència de Solvay a Brussel·les, Albert Einstein i Niels Bohr van protagonitzar un dels debats científics més importants de la història sobre la mecànica quàntica. Després d'una sessió esgotadora on Einstein va pronunciar la seva mítica frase "Déu no juga als daus", va decidir anar directament cap a la cafeteria central per esmorzar, sortint de l'hotel amb la seva bicicleta a una velocitat constant de $\mathbf{18\ km/h}$. En Bohr, que estava a l'estació de trens situada a \textbf{320 decàmetres} en línia recta de l'hotel d'Einstein, s'adona que l'alemany s'ha deixat els seus apunts. Bohr surt corrent per atrapar-lo a una velocitat de \textbf{250 m/minut} en direcció a la cafeteria (cap a Einstein). \textbf{Atenció:} Com que Bohr ha hagut de recollir els papers, surt exactament \textbf{2 minuts} més tard que Einstein.

\begin{enumerate}[label=\textbf{\alph*)}]
  \item Escriu les equacions de posició $x(t)$ per a tots dos físics.
  \item Calcula en quin instant de temps (en minuts des de la sortida d'Einstein) es trobaran.
  \item A quina distància (en quilòmetres) de l'hotel d'Einstein estaran quan es creuin?
  \item Dibuixa les 3 gràfiques del moviment ($x-t$, $v-t$, $a-t$) on hi apareguin, en els mateixos eixos, les línies de tots dos cossos.
\end{enumerate}

\drawplots{x}
\gradebox{2,5}

\newpage

% ================================================================
% PÀGINA 3: EXERCICI 2
% ================================================================
\exercici{2,5}{MRUA · ThrustSSC: Trencant la barrera del so}

El 15 d'octubre de 1997, al desert de Black Rock a Nevada, el pilot britànic Andy Green va fer història a bord del ThrustSSC, convertint-se en el primer vehicle terrestre en trencar la barrera del so de manera oficial. Durant una prova prèvia, el ThrustSSC arrenca des del repòs i accelera constantment fins a assolir els $\mathbf{1080\ km/h}$ en un temps de \textbf{0,4 minuts}. Just en assolir aquesta velocitat punta, salta una alarma al quadre de comandaments. L'Andy ha d'activar el paracaigudes de frenada d'emergència, aturant el cotxe completament i de manera uniforme en un espai de \textbf{45\,000 decímetres}.

\begin{enumerate}[label=\textbf{\alph*)}]
  \item Quina ha estat l'acceleració (en m/s²) en cadascun dels dos trams?
  \item Quina distància total (en quilòmetres) ha recorregut des que ha arrencat fins que s'ha aturat definitivament?
  \item Dibuixa les 3 gràfiques ($x-t$, $v-t$, $a-$t) de tot el moviment complet (unint els dos trams).
\end{enumerate}

\drawplots{x}
\gradebox{2,5}

\newpage

% ================================================================
% PÀGINA 4: EXERCICI 3
% ================================================================
\exercici{2,5}{Caiguda Lliure · Galileu i la Torre de Pisa}

L'any 1589, el geni italià Galileo Galilei va voler demostrar al món que la teoria de la caiguda dels cossos d'Aristòtil estava totalment equivocada. La llegenda explica que va pujar a dalt de tot de la famosa Torre Inclinada de Pisa. Galileu deixa caure des del balcó principal, situat a \textbf{5,4 decàmetres} d'altura respecte a terra, una pesant bola de canó de ferro pur ($v_0 = 0$). Exactament en el mateix instant, el seu fidel ajudant Vincenzo, situat a una finestra més baixa a \textbf{35\,000 mil·límetres} d'altura, llança una altra bola verticalment cap avall amb una velocitat inicial de \textbf{800 cm/s}.

\begin{enumerate}[label=\textbf{\alph*)}]
  \item Arribaran a xocar a l'aire les dues boles abans de tocar el terra? Si és que sí, calcula l'altura (en metres) i l'instant de temps exacte de l'impacte.
  \item Si cap de les dues boles xoca amb l'altra, quin dels dos objectes impactarà primer contra la plaça i amb quina velocitat (en km/h) ho farà cadascun?
  \item Dibuixa les 3 gràfiques ($y$-$t$, $v$-$t$, $a$-$t$) d'ambdós objectes.
\end{enumerate}

\drawplots{y}
\gradebox{2,5}

\newpage

% ================================================================
% PÀGINA 5: EXERCICI 4
% ================================================================
\exercici{2,5}{Tir Vertical · Wimbledon 2019: Federer vs Djokovic}

Som al 14 de juliol de 2019, a la llegendària pista central de Wimbledon. És la gran final de Grand Slam més llarga de la història. Roger Federer es prepara per executar el seu servei en un punt crític. Llança la pilota verticalment cap amunt des d'una altura inicial de \textbf{165 centímetres} respecte a la gespa, imprimint-li una velocitat inicial de $\mathbf{18\ km/h}$. El punt òptim per impactar la pilota és exactament al seu punt mort superior ($y_{\max}$), on la velocitat és momentàniament zero. Malauradament, un reflex del sol l'encega, no arrive a tocar la pilota, i la deixa caure lliurement fins a terra.

\begin{enumerate}[label=\textbf{\alph*)}]
  \item A quina altura màxima, respecte al terra de gespa, tenia pensat en Federer colpejar la pilota ($y_{\max}$)?
  \item Quant de temps (en minuts) va estar volant en total la pilota, des que va sortir de la mà del suís fins que va tocar finalment el terra de la pista per primer cop?
  \item Dibuixa les 3 gràfiques ($y$-$t$, $v$-$t$, $a$-$t$) de l'acció completa.
\end{enumerate}

\drawplots{y}
\gradebox{2,5}

\newpage

% ================================================================
% PÀGINA 6: DEMOSTRACIONS I EXTRA
% ================================================================
\section*{Punts Extra: Demostracions Teòriques}

\textbf{Instruccions:} Per a cadascuna de les següents deduccions, primer explica sota quines condicions són aplicables i després desenvolupa el procés matemàtic pas a pas.

\vspace{10pt}

\section*{Exercici Extra 1 \textnormal{(2 punts)} \hfill \normalsize\normalfont\textit{Deducció de fórmules}}

A partir de les equacions generals del MRUA per a un llançament vertical cap amunt:

\begin{enumerate}[label=\textbf{\alph*)}, itemsep=15pt]
  \item \textbf{Temps Màxim ($t_{\max}$):} Demostra que el temps necessari per arribar al punt més alt del llançament es pot calcular com $t_{\max} = \frac{v_0}{g}$. \textit{(1 punt)}
  \item \textbf{Altura Màxima ($y_{\max}$):} Utilitzant el resultat anterior, demostra que l'altura màxima assolida (partint d'una altura inicial $\mathbf{y_0 = h}$) és $y_{\max} = h + \frac{v_0^2}{2g}$. \textit{(1 punt)}
\end{enumerate}

\vfill

\noindent
\begin{tabularx}{\textwidth}{|>{\raggedright\arraybackslash}X|>{\centering\arraybackslash}p{5cm}|}
  \hline
  \rule{0pt}{22pt}\textbf{Demostració A ($t_{\max}$)} & \Large\hspace{2.5cm}~/~1,0 \\[8pt]
  \hline
  \rule{0pt}{22pt}\textbf{Demostració B ($y_{\max}$)} & \Large\hspace{2.5cm}~/~1,0 \\[8pt]
  \hline
  \rule{0pt}{22pt}\textbf{\Large TOTAL EXTRA} & \Large\textbf{\hspace{2.5cm}~/~2,0} \\[8pt]
  \hline
\end{tabularx}

\end{document}